\documentclass[10pt,a4paper]{amsart}
\usepackage[margin=19mm]{geometry}
\usepackage{amsmath,amssymb,mathtools}
\usepackage[scr=rsfs,cal=euler]{mathalfa}
\usepackage{xfrac}
\usepackage[unicode,hidelinks,bookmarks=false]{hyperref}
\newcommand{\ie}{i.e.\ }
\newcommand{\cf}{cf.\ }
\newcommand{\resp}{respectively\ }
\newcommand{\Subsection}{\subsection}
\providecommand{\nobreakdash}{\nobreak\hbox{-}}
\setlength{\emergencystretch}{2em}
\multlinegap0pt
\usepackage{amssymb}
\usepackage{tikz}
\newtheorem{lemma}{Lemma}
\title{On the gauge invariance of the Kuperberg invariant of certain high genus framed 3-manifolds}
\author{Liang Chang, Yilong Wang, Saifei Zhai}
\date{Source: arXiv:2601.19485v1 (CC BY 4.0)}
\begin{document}
\maketitle

\noindent\emph{Source and licence.} On the gauge invariance of the Kuperberg invariant of certain high genus framed 3-manifolds. Version arXiv:2601.19485v1 (CC BY 4.0), distributed by arXiv under the Creative Commons Attribution 4.0 International licence, http://creativecommons.org/licenses/by/4.0/, as recorded in the arXiv licence metadata for this exact version and shown on the abstract page https://arxiv.org/abs/2601.19485v1. Full licence notice: \texttt{licenses/arxiv-2601.19485-CC-BY-4.0.txt}.\par\smallskip

\noindent\textit{Editorial scope.} Counted unit: the article's lemma lem:3 (source0.tex lines 636--644). The article's complete original proof of it is delivered with it (source0.tex lines 646--681, one \texttt{proof} environment of 36 lines). No mathematics is written, repaired or re-derived: the statement and the proof are the article's own lines, in the article's own notation, and no retained line is substituted or re-flowed.  Retained uncounted as the source-local context the proof needs: lines 505--521; lines 614--619.  Nothing was omitted from any retained span, so the package carries no omission record.  No retained line contains a citation, so the wrapper carries no bibliography and no bibliography entry.  No retained line contains a figure environment, an image-inclusion command or drawing code, so no image is delivered and the package declares no asset dependency.  Editorial formatting only, in addition: an AMS \texttt{amsart} preamble that restates the article's own declarations and macros used by the retained lines, namely the environments \texttt{lemma} on separate counters, together with the standard packages those lines need.  The retained environments print in this document as Lemma 1 in order of appearance, whereas the article prints the same items with the numbering of its own full document.  The complete native source of the article as distributed in the arXiv e-print for this exact version is bundled unchanged as \texttt{sources/193458/source0.tex}.
\begin{lemma}\label{lemma1}
For any $x\in H$ and $Y, Z\in \mathrm{End}(H)$, \\
$(1)$ \[\begin{aligned}
&\lambda S (S^{-1}(\Lambda^2_{(1)})Y(\Lambda^2_{(7)},\Lambda^1_{(5)},\Lambda^1_{(2)},\Lambda^2_{(8)},\Lambda^2_{(5)},\Lambda^1_{(7)},\Lambda^2_{(3)})\Lambda^1_{(9)})\\
&\lambda S(xZ(\Lambda^1_{(3)},\Lambda^{1}_{(6)},\Lambda^2_{(4)},\Lambda^1_{(8)},\Lambda^2_{(2)},\Lambda^1_{(1)},\Lambda^1_{(4)},\Lambda^2_{(6)})\Lambda^2_{(9)})\\
=&\lambda S (S^{-1}(\Lambda^2_{(1)})Y(\Lambda^2_{(7)}S(x_{(2)}),\Lambda^1_{(5)}S(x_{(11)}),\Lambda^1_{(2)}S(x_{(14)}),\Lambda^2_{(8)}S(x_{(1)}),\Lambda^2_{(5)}S(x_{(4)}),\Lambda^1_{(7)}S(x_{(9)}),\\
&\Lambda^2_{(3)}S(x_{(6)}))\Lambda^1_{(9)})\\
&\lambda S(Z(\Lambda^1_{(3)}S(x_{(13)}),\Lambda^{1}_{(6)}S(x_{(10)}),\Lambda^2_{(4)}S(x_{(5)}),\Lambda^1_{(8)}S(x_{(8)}),\Lambda^2_{(2)}S(x_{(7)}),\Lambda^1_{(1)}S(x_{(15)}),\Lambda^1_{(4)}\\
&S(x_{(12)}),\Lambda^2_{(6)}S(x_{(3)}))\Lambda^2_{(9)})\\
\end{aligned}\]
$(2)$ \[\begin{aligned}
	&\lambda S (Y(\Lambda^2_{(1)},\Lambda^2_{(7)},\Lambda^1_{(5)},\Lambda^1_{(2)},\Lambda^2_{(8)},\Lambda^2_{(5)},\Lambda^1_{(7)})S^{-1}(\Lambda^2_{(3)})\Lambda^1_{(9)})\\
	&\lambda S(Y(\Lambda^1_{(3)},\Lambda^{1}_{(6)},\Lambda^2_{(4)},\Lambda^1_{(8)},\Lambda^2_{(2)},\Lambda^1_{(1)},\Lambda^1_{(4)},\Lambda^2_{(6)})S^{-1}(x)\Lambda^2_{(9)})\\
	=&\lambda S (Y(x_{(1)}\Lambda^2_{(1)},x_{(14)}\Lambda^2_{(7)},x_{(7)}\Lambda^1_{(5)},x_{(4)}\Lambda^1_{(2)},x_{(15)}\Lambda^2_{(8)},x_{(12)}\Lambda^2_{(5)},x_{(9)}\Lambda^1_{(7)}S^{-1}(\Lambda^2_{(3)})\Lambda^1_{(9)})\\
	&\lambda S(Z(x_{(5)}\Lambda^1_{(3)},x_{(8)}\Lambda^{1}_{(6)},x_{(11)}\Lambda^2_{(4)},x_{(10)}\Lambda^1_{(8)},x_{(2)}\Lambda^2_{(2)},x_{(3)}\Lambda^1_{(1)},x_{(6)}\Lambda^1_{(4)},x_{(13)}\Lambda^2_{(6)})\Lambda^2_{(9)})\\
\end{aligned}\]
\end{lemma}

\begin{lemma}\label{lemma2}
	For any $x\in H$ and $X, Y\in \mathrm{End}(H)$,
\[\begin{aligned}
	\lambda S(X(\Lambda_{(1)}x)Y(\Lambda_{(2)})\Lambda_{(3)})=\lambda S(S^{-1}(x_{(1)})X(\Lambda_{(1)})Y(\Lambda_{(2)}S(x_{(2)}))\Lambda_{(3)}).
\end{aligned}\]
\end{lemma}

\begin{lemma}\label{lem:3}
For any $x, w_1, w_2, w_3\in H$,
\[\begin{aligned}
	&\lambda S (S^{-1}(\Lambda^2_{(1)})S^{-1}(\Lambda^2_{(7)})S^{-3}(\Lambda^1_{(5)})w_2S^{-1}(\Lambda^1_{(2)}x_{(2)})S^{-1}(\Lambda^2_{(8)})S^{-3}(\Lambda^2_{(5)})S^{-2}(\Lambda^1_{(7)})S^{-1}(\Lambda^2_{(3)})\Lambda^1_{(9)})\\
	&\lambda S(S^{-2}(\Lambda^1_{(3)}x_{(3)})w_3S^{-4}(\Lambda^1_{(6)})S^{-3}(\Lambda^2_{(4)})S^{-2}(\Lambda^1_{(8)})S^{-1}(\Lambda^2_{(2)})S^{-2}(\Lambda^1_{(1)}x_{(1)})w_1S^{-4}(\Lambda^1_{(4)})S^{-2}(\Lambda^2_{(6)})\Lambda^2_{(9)})\\
	=&\lambda S (S^{-1}(\Lambda^2_{(1)})S^{-1}(\Lambda^2_{(7)})S^{-3}(x_{(2)}\Lambda^1_{(5)})w_2S^{-1}(\Lambda^1_{(2)})S^{-1}(\Lambda^2_{(8)})S^{-3}(\Lambda^2_{(5)})S^{-2}(\Lambda^1_{(7)})S^{-1}(\Lambda^2_{(3)})\Lambda^1_{(9)})\\
	&\lambda S(S^{-2}(\Lambda^1_{(3)})w_3S^{-4}(x_{(3)}\Lambda^1_{(6)})S^{-3}(\Lambda^2_{(4)})S^{-2}(\Lambda^1_{(8)})S^{-1}(\Lambda^2_{(2)})S^{-2}(\Lambda^1_{(1)})w_1S^{-4}(x_{(1)}\Lambda^1_{(4)})S^{-2}(\Lambda^2_{(6)})\Lambda^2_{(9)})\\
\end{aligned}\]
\end{lemma}

\begin{proof} 
Apply Lemma \ref{lemma2} for $\Lambda^1_{(1)}x_{(1)}\otimes\Lambda^1_{(2)}x_{(2)}\otimes\Lambda^1_{(3)}x_{(3)}$. The left hand side becomes
\[
\begin{aligned}
	L=&\lambda S (S^{-1}(x_{(1)})S^{-1}(\Lambda^2_{(1)})S^{-1}(\Lambda^2_{(7)})S^{-2}(x_{(5)})S^{-3}(\Lambda^1_{(5)})w_2S^{-1}(\Lambda^1_{(2)})S^{-1}(\Lambda^2_{(8)})S^{-3}(\Lambda^2_{(5)})S^{-2}(\Lambda^1_{(7)})\\
	&S^{-1}(x_{(3)})S^{-1}(\Lambda^2_{(3)})\Lambda^1_{(9)})\\
	&\lambda S(S^{-2}(\Lambda^1_{(3)})w_3S^{-4}(\Lambda^1_{(6)})S^{-2}(x_{(4)})S^{-3}(\Lambda^2_{(4)})S^{-2}(\Lambda^1_{(8)})S^{-1}(x_{(2)})S^{-1}(\Lambda^2_{(2)})S^{-2}(\Lambda^1_{(1)})w_1S^{-4}(\Lambda^1_{(4)})\\
	&S^{-3}(x_{(6)})S^{-2}(\Lambda^2_{(6)})\Lambda^2_{(9)})\\
\end{aligned}
\]
Similarly, Lemma \ref{lemma2} is applied to $S^{-1}(x_{(1)})S^{-1}(\Lambda^2_{(1)})=S^{-1}(\Lambda^2_{(1)}x_{(1)})$. Then we get that
\[
\begin{aligned}
L=&\lambda S (S^{-1}(\Lambda^2_{(1)})x_{(3)}S^{-1}(\Lambda^2_{(7)})S^{-2}(x_{(12)})S^{-3}(\Lambda^1_{(5)})w_2S^{-1}(\Lambda^1_{(2)})x_{(2)}S^{-1}(\Lambda^2_{(8)})S^{-2}(x_{(5)})S^{-3}(\Lambda^2_{(5)})\\
&S^{-2}(\Lambda^1_{(7)})S^{-1}(x_{(10)})x_{(7)}S^{-1}(\Lambda^2_{(3)})\Lambda^1_{(9)})\\
&\lambda S(S^{-1}(x_{(1)})S^{-2}(\Lambda^1_{(3)})w_3S^{-4}(\Lambda^1_{(6)})S^{-3}(x_{(11)})S^{-2}(x_{(6)})S^{-3}(\Lambda^2_{(4)})S^{-2}(\Lambda^1_{(8)})S^{-1}(x_{(9)})x_{(8)}\\
&S^{-1}(\Lambda^2_{(2)})S^{-2}(\Lambda^1_{(1)})w_1S^{-4}(\Lambda^1_{(4)})S^{-3}(x_{(13)})S^{-2}(\Lambda^2_{(6)})S^{-1}(x_{(4)})\Lambda^2_{(9)})\\
=&\lambda S (S^{-1}(\Lambda^2_{(1)})x_{(3)}S^{-1}(\Lambda^2_{(7)})S^{-2}(x_{(6)})S^{-3}(\Lambda^1_{(5)})w_2S^{-1}(\Lambda^1_{(2)})x_{(2)}S^{-1}(\Lambda^2_{(8)})S^{-2}(x_{(5)})S^{-3}(\Lambda^2_{(5)})\\
&S^{-2}(\Lambda^1_{(7)})S^{-1}(\Lambda^2_{(3)})\Lambda^1_{(9)})\\
&\lambda S(S^{-1}(x_{(1)})S^{-2}(\Lambda^1_{(3)})w_3S^{-4}(\Lambda^1_{(6)})S^{-3}(\Lambda^2_{(4)})S^{-2}(\Lambda^1_{(8)})S^{-1}(\Lambda^2_{(2)})S^{-2}(\Lambda^1_{(1)})w_1S^{-4}(\Lambda^1_{(4)})S^{-3}(x_{(7)})\\
&S^{-2}(\Lambda^2_{(6)})S^{-1}(x_{(4)})\Lambda^2_{(9)})\\
\end{aligned}
\]
By Lemma \ref{lemma1}(2), we get that
\[
\begin{aligned}
L=&\lambda S (S^{-1}(\Lambda^2_{(1)})S^{-1}(x_{(4)})x_{(3)}S^{-1}(\Lambda^2_{(7)})S^{-1}(x_{(17)})S^{-2}(x_{(20)})S^{-3}(\Lambda^1_{(5)})S^{-3}(x_{(10)})w_2S^{-1}(\Lambda^1_{(2)})\\
&S^{-1}(x_{(7)})x_{(2)}S^{-1}(\Lambda^2_{(8)})S^{-1}(x_{(18)})S^{-2}(x_{(19)})S^{-3}(\Lambda^2_{(5)})S^{-3}(x_{(15)})S^{-2}(x_{(12)})S^{-2}(\Lambda^1_{(7)})S^{-1}(\Lambda^2_{(3)})\Lambda^1_{(9)})\\
&\lambda S(S^{-1}(x_{(1)})S^{-2}(x_{(8)})S^{-2}(\Lambda^1_{(3)})w_3S^{-4}(x_{(11)})S^{-4}(\Lambda^1_{(6)})S^{-3}(\Lambda^2_{(4)})S^{-3}(x_{(14)})S^{-2}(x_{(13)})S^{-2}(\Lambda^1_{(8)})\\
&S^{-1}(\Lambda^2_{(2)})S^{-1}(x_{(5)})S^{-2}(x_{(6)})S^{-2}(\Lambda^1_{(1)})w_1S^{-4}(x_{(9)})S^{-4}(\Lambda^1_{(4)})S^{-3}(x_{(21)})S^{-2}(x_{(16)})S^{-2}(\Lambda^2_{(6)})\Lambda^2_{(9)})\\
=&\lambda S (S^{-1}(\Lambda^2_{(1)})S^{-1}(\Lambda^2_{(7)})S^{-3}(\Lambda^1_{(5)})S^{-3}(x_{(2)})w_2S^{-1}(\Lambda^1_{(2)})S^{-1}(\Lambda^2_{(8)})S^{-3}(\Lambda^2_{(5)})S^{-2}(\Lambda^1_{(7)})S^{-1}(\Lambda^2_{(3)})\Lambda^1_{(9)})\\
&\lambda S(S^{-2}(\Lambda^1_{(3)})w_3S^{-4}(x_{(3)})S^{-4}(\Lambda^1_{(6)})S^{-3}(\Lambda^2_{(4)})S^{-2}(\Lambda^1_{(8)})S^{-1}(\Lambda^2_{(2)})S^{-2}(\Lambda^1_{(1)})w_1S^{-4}(x_{(1)})S^{-4}(\Lambda^1_{(4)})\\
&S^{-2}(\Lambda^2_{(6)})\Lambda^2_{(9)})
\end{aligned}\]
This completes the proof.
\end{proof}

\end{document}
